\begin{result}{F.3.36}{The characteristic-zero comparison on all abelian tensor motives}{Known source theorem; category and descent maps proved}
Let \(k\subset\C\) be a field. For rational algebraic cycles on every
smooth projective \(X/k\) of dimension at most four, and on every
abelian variety \(A/k\) of arbitrary dimension, Betti homological
equivalence and numerical equivalence coincide. The same assertion
holds for geometric \(\ell\)-adic homological equivalence for every
prime \(\ell\). Consequently, the identity on cycles induces a tensor
equivalence between the additive, idempotent-complete categories of
homological and numerical motives generated by all \(h(A)(m)\),
where \(A/k\) is abelian and \(m\in\Z\). The equivalence acts as
the identity on the named generators and on their correspondence
classes. This assertion includes arbitrary products, tensor powers,
duals and numerical direct summands of those generators.
\end{result}
\begin{proof}
The all-dimensional complex abelian-variety theorem follows from
the algebraicity of the inverse Lefschetz operator in
\cite[Theorem 2A11, p.377]{AuditKleiman1968} and Hodge positivity;
the numerical--homological conclusion is stated explicitly in the
same source, Remark~3.10 on p.380. That remark also proves the
assertion for every smooth projective complex variety of dimension
at most four, using the Lefschetz \((1,1)\) theorem and Hodge
positivity. We use those stated theorems.

Here is the field-extension step for the actual cycle groups over
\(k\). Let \(z\in\operatorname{CH}^r(X)_{\Q}\) be numerically
zero over \(k\), where \(X\) is any of the stated varieties.
For a finite extension
\(L/k\), projection \(p:X_L\to X\), and a complementary cycle
\(w\) on \(X_L\), the projection formula and the degree of a
zero-cycle over its field of definition give
\[
 [L:k]\deg_L(z_L\cdot w)
   =\deg_k(z\cdot p_*w)=0.                     \tag{F.3.36a}
\]
Now let \(w_{\C}\) be any complementary cycle on \(X_{\C}\).
It suffices by linearity to treat an integral subvariety. Spread its
defining equations over an integral finite-type \(k\)-scheme \(T\),
and shrink \(T\) so that the resulting family
\(\mathcal W\subset X\times_kT\) is flat of relative dimension
\(r\), and \(T\) is smooth. The given complex cycle is the base
change of the generic fibre. The relative intersection of the
constant cycle \(z\) with \(\mathcal W\), pushed forward to \(T\),
is \(d[T]\) in \(\operatorname{CH}^0(T)_{\Q}\). Compatibility
with refined pullback to a point identifies \(d\) with the
intersection degree on each fibre. Choose a closed point \(t\in T\);
its residue field \(L/k\) is finite. Equation~(F.3.36a) makes
the fibre degree zero, so \(d=0\), and the degree on the complex
fibre is zero as well. Thus \(z_{\C}\) is numerically zero.
The cited complex theorem gives
\(\operatorname{cl}_B(z_{\C})=0\). Conversely, a zero Betti class
has zero intersection degree with every complementary cycle over
\(k\), by compatibility of the cycle class with products and trace.
This proves the equality over \(k\).

The smooth proper comparison isomorphism
\[
 H_B^{2r}(X(\C),\Q)\otimes_{\Q}\Q_\ell
   \simeq H^{2r}_{\acute e t}(X_{\overline k},\Q_\ell)
\]
(with the corresponding Tate twists) carries the two cycle classes
to each other. Faithful scalar extension therefore gives exactly
the same equivalence relation for geometric \(\ell\)-adic classes.

To identify every morphism in the stated categories, use the
contravariant convention of (F.3.35c). For arbitrary abelian
\(A,B\) and integers \(m,n\), the quotient map is the identity on
the cycles in
\[
 \operatorname{CH}^{\dim A+n-m}(A\times B)_{\Q},
                                                        \tag{F.3.36b}
\]
and its homological and numerical kernels have just been proved
equal, because \(A\times B\) is again an abelian variety. It is
therefore an isomorphism on these Hom spaces. The composition
formula uses the same pullback, intersection product and proper
pushforward on both sides, so these isomorphisms form a functor.
They remain isomorphisms on matrix Hom spaces for finite sums and
on \(f\operatorname{Hom}e\) for every pair of idempotents \(e,f\).
Every numerical idempotent lifts uniquely through the endomorphism
isomorphism, and injectivity makes the lift idempotent. Hence the
functor is fully faithful and essentially surjective onto the
specified numerical category.

Finally, the actual product correspondence gives
\(h(A)(m)\otimes h(B)(n)=h(A\times B)(m+n)\), and Poincar\'e
duality gives \(h(A)^{\vee}=h(A)(\dim A)\). Thus all tensor
products and duals remain among the generators, and the preceding
argument already covers all their projectors. Exterior products
of correspondences are unchanged by the quotient, which supplies
the tensor compatibility. The all-smooth-projective numerical
quotient in F.3.5 still requires its own assertion on products
outside this abelian category.
\end{proof}
